\chapter{Implementation}
\label{chap:implementation}
\textit{This chapter describes how \sysname{} is implemented: the technology stack, the agent workflow on LangGraph, the backend, the web app, the testing and continuous integration process, the deployment, and a demonstration of the prototype on a real idea.}

\section{Technology Stack}
\label{sec:impl-stack}

Table~\ref{tab:impl-stack} lists the technologies used for each concern of the system, from the orchestration of the agents to the tooling of the team.

\begin{table}[H]
    \centering
    \caption{Technology stack.}
    \label{tab:impl-stack}
    \begin{tabularx}{\linewidth}{|l|X|}
        \hline
        \textbf{Concern} & \textbf{Choice} \\
        \hline
        Agent orchestration & LangGraph, LangChain Core \\
        LLM access & LangChain chat models with provider packages (OpenAI, Anthropic, Google, Ollama); OpenAI-compatible gateways \\
        Research search & Tavily web search, or a keyless chain of DuckDuckGo and Wikipedia; OpenAlex and arXiv for paper abstracts \\
        Backend & Python 3.12, FastAPI, Uvicorn, Pydantic v2, structlog \\
        Persistence & SQLAlchemy 2 and Alembic migrations on SQLite (aiosqlite). PostgreSQL with pgvector is the production target \\
        Frontend & Next.js 16 (App Router), React 19, TypeScript (strict), Tailwind CSS v4, shadcn/ui, TanStack Query, Zustand, React Hook Form with Zod \\
        API client & Generated from the OpenAPI schema of the backend \\
        Quality & Ruff, mypy (strict), import-linter, pytest; ESLint, Prettier, Vitest with Testing Library and MSW, Playwright \\
        Tooling & uv, pnpm, Docker Compose, GitHub Actions \\
        \hline
    \end{tabularx}
\end{table}

The main choices are motivated as follows:
\begin{itemize}
    \item \textbf{LangGraph} represents the workflow as an explicit state graph with conditional edges and a checkpointer (Section~\ref{sec:bg-langgraph}). The conditional edges carry the loops and the graph can be run and tested without the API; the backend saves the state after every step, so pausing, resuming and retrying from a failed step are its own code, not a service of the runtime.
    \item A \textbf{modular monolith} keeps one deployable backend while enforcing module boundaries with import rules. For a team of three, it avoids the operational cost of services while keeping the agents independent of the web layer, which the evaluation harness requires.
    \item A \textbf{relational database} holds the users, runs, events, checkpoints and the trace graph in one store (Section~\ref{sec:sd-database}). SQLite keeps the development setup to one file; the models and migrations do not depend on it.
    \item \textbf{Generated API types} make the web app and the backend agree on every request and response at compile time.
\end{itemize}

\section{Agent Workflow Implementation}
\label{sec:impl-agents}

\subsection{State and Graph}
\label{sec:impl-graph}

The workflow has a shared \emph{state} that every step reads and updates. It holds the inputs of the run (the project, the initial idea and the settings of the run), references to the artifacts and items the run has produced, at the version each step used, the current stage, the status of the run, and the few values that decide the next step: the answers given at the checkpoints, the verdicts of the analyst and of the validator, and the step chosen by the intent router. The content of the artifacts is kept in the database, not in the state, so the state stays small and every saved state can still be read later.

The workflow of version~2 has twenty-two steps connected by 36 transitions. The research steps follow each other in a fixed order (planner, corpus builder, stakeholder miner, persona generator); every other transition is conditional and is decided from the state alone, for example:
\begin{itemize}
    \item after the interviewer: to the persona agent for a new turn, or to the analyst when all interviews are complete;
    \item after the analyst: to the rework control node when the analysis needs interviews, otherwise to the analysis checkpoint;
    \item after a checkpoint: along the edge of the option the user chose;
    \item after the validator: to the end on a pass or when the maximum number of review rounds is reached, otherwise back to the requirement engineer.
\end{itemize}
A checkpoint does not keep the run waiting in memory. The run stops with the status \emph{interrupted} and a review request is stored; when the user answers, the run starts again from the saved state, at the step that the chosen option leads to. A run that waits for the user, even for days, therefore holds no connection and no running process.

The steps and transitions are defined once, in the workflow itself. The description of the workflow that the web app draws (labels, stages, agent roles and checkpoint options) is derived from it, and an automated test checks that the two agree.

\subsection{Prompts}
\label{sec:impl-prompts}

The prompt of every agent role is a text document with a version number, kept and versioned together with the workflow (Appendix~\ref{apx:prompts}). A short header names the role, the version, the expected output and the inputs. The body has a \emph{System} part, with the required sections Role, Goal, Inputs, Procedure and Rules and the optional Output and Examples, and a \emph{User} part with the sections Context and Task. The system part is fixed for every call of a role, so that it is served from the prompt cache of the provider; the user part carries the changing data, each input enclosed in a named tag (for example, a tag that encloses the interview transcripts), and every prompt states that text inside the tags is data, not instructions. The structure of every prompt is checked automatically, and the inputs are filled in at each call. The version of each role comes from the workflow definition, so one role can move to a new prompt version, as the persona generator did for version~2, while the others keep theirs.

\subsection{Model Layer and Structured Output}
\label{sec:impl-llm}

Agents never name a model: they ask for a model by role, and the model access module chooses the provider and the model of the role from a configuration file, the settings of the user and the settings of the run, then fixes this choice for the whole run. Each call asks the model for structured output that must conform to the output schema of the role. A usage gate around every call checks the cost limit of the run and the budget of the user before the call and records the token usage, including prompt cache reads and writes, after it. A response that cannot be parsed, for example a streamed generation that degenerates into a truncated tool call, is retried twice before the step fails.

After the call, the step turns the output into items and links and saves them; the integrity checks of the domain rules reject an inconsistent output. Some nodes add checks of their own, such as the corpus builder, which rejects a document whose URL the search did not return.

\section{Backend Implementation}
\label{sec:impl-backend}

The backend follows the modules of Section~\ref{sec:sd-backend-modules}. Each feature module (accounts, projects, conversations, runs, artifacts, documents, exports and the assistant) contains its API endpoints, the formats of its requests and responses, its business logic and its database tables, built with FastAPI, Pydantic and SQLAlchemy; the database schema is created by versioned migrations.

\paragraph{Run execution.} Posting the first message of a thread creates a run with the status \emph{queued} and returns at once. The run is then executed in the background: the backend fixes the settings of the run, gives the workflow access to the database, the document search, the project documents and the cost control, and runs the workflow from its last saved state. For every step it records which step ran, the model used, the tokens and cost, and the items produced. On startup, the backend recovers the runs left by a stopped process as described in Section~\ref{sec:sd-failures}.

\paragraph{Event log and streaming.} Every event of a run (a token, a step started or finished, an artifact changed, a review requested, the run finished) is appended to an event log in the database with an increasing sequence number. The streaming endpoint reads the log and sends the events as Server-Sent Events; a client that reconnects passes the last sequence number it received and gets the events after it.

\paragraph{Assistant.} Every posted message is handed to the assistant as a job of the server (Section~\ref{sec:pa-assistance}). The assistant calls its model with structured output that either answers the message, asks for one more read through one of its tools, or sends the message on to the requirements workflow. Its tools only read the database of the project. The commands of the answer are kept only when they are listed in the catalog of commands and name a document or an item of the project, and the usage of the call is recorded so that the budget of the user includes it.

\paragraph{Artifacts and traceability.} The artifact service validates every write with the domain checks and stores a new revision. The traceability endpoints compute the trace of an item, the traceability matrix and the coverage of a project from the trace links. Exports render the specification, the backlog, the matrix and the transcripts into files.

\section{Web Application}
\label{sec:impl-frontend}

The web app follows the layers of Section~\ref{sec:sd-frontend}. Its main screens are the project list, the project page with its threads and documents, and the thread page, which combines the chat, the workflow graph with the progress of the current run, the review of a pending checkpoint, and an artifact panel. The graph, its stages and the review forms are generated from the workflow definition: a checkpoint's options become buttons and its questions become form fields, with the validation rules of the definition. The artifact panel renders each artifact kind from its content schema, with editors for the items and the links. Every agent is shown with a small animated face. Its mood follows the state of the work: working while the run runs, waiting at a review, glad when the run succeeds, worried when it fails and asleep when it is cancelled, among eight moods in all. Beside the face, an \emph{act} shows what the agent is doing, such as searching, reading, interviewing or writing; the seventeen acts form a catalog, and every stage and tool of the workflow definition names the acts it performs, so that a new stage needs no change of the web app. The prompt input also accepts dictation, and answers can be read aloud. In mock mode, simulated responses replace the API inside the browser, which is used both for development and for the end-to-end tests. Screenshots of the main screens are shown in Section~\ref{sec:sd-ui}.

\section{Testing and Continuous Integration}
\label{sec:impl-testing}

\paragraph{Backend.} The backend has 759 automated tests, of which 10 are currently skipped; all the others pass. They cover the domain rules without a database; the consistency of the workflow description, the item formats and the prompts; the workflow itself, run with simulated language models; the document search; the automatic decisions; and the API end to end, including authentication, editing, exports, the event stream with review requests, and the execution of runs. Static checks verify code style and types, and enforce the allowed dependencies between the modules.

\paragraph{Frontend.} The web app has 615 files of unit and component tests, and 22 suites of end-to-end tests that drive the application in a browser. On every change, the continuous integration of the web app checks types and code style, runs the unit tests and reports their coverage, builds the application and runs the end-to-end tests; a second job captures screenshots of the screens.

The continuous integration of the backend runs on every change as well: it checks the code style, the types and the allowed dependencies between the modules, runs the automated tests, and runs end-to-end tests of the web application against the real backend.

\section{Deployment}
\label{sec:impl-deployment}

For development, the API and the web app run in containers that reload when their source changes, with the database stored in a single file. Secrets, such as the keys of the model providers, are supplied by the operator and are never stored with the source. The models available to the agents are listed in a configuration that can point to the public providers or to a compatible gateway, and a script creates a demo user with sample projects. The production layout of Section~\ref{sec:sd-deployment} is the target of the next phase.

\section{Prototype Demonstration}
\label{sec:impl-demo}
\label{sec:eval-case-study}

To check that the prototype works end to end outside the tests, the system was run once through its API, with real LLM calls and real search, on the idea ``\textit{Build a platform that helps university students find career mentors}''. The project was created with Vietnamese as output language, and the members uploaded one short document of notes collected at a university. The run took about 33 minutes and went from the idea to the interviews. It shows that the first half of the workflow, from planning to persona-based interviews, runs on real models and real documents and keeps its evidence chain; it is a demonstration of the prototype, not a measurement of the quality of the requirements, which the evaluation planned in Chapter~\ref{chap:next-phase} addresses. The paragraphs below give an overview of what each step produced.

\paragraph{Plan.} The planner listed six stakeholder types: students seeking career guidance, potential mentors, career services staff, university administrators, platform administrators, and employers and industry partners. Its topics include the current channels students use and their frustrations, the outcomes expected from mentoring, matching criteria, interaction formats, and credential verification.

\paragraph{Corpus and observations.} The corpus builder kept three documents: the uploaded notes, a journal article on the impact of mentoring in higher education on student careers, and an article on developing strategic partnerships between companies and universities. The stakeholder miner extracted eight observations, each with a verbatim quote, for example:
\begin{itemize}
    \item students, \emph{advisor matching by field and free time}: ``Tim co van theo nganh nghe va gio ranh'' (find an advisor by field and free time);
    \item students, \emph{quick booking in a few clicks}: ``Book lich nhanh trong vai cu bam'' (book a session in a few clicks);
    \item advisors, \emph{filtering requests before accepting}: ``Loc yeu cau truoc khi nhan'' (filter requests before accepting them);
    \item companies, \emph{difficulty establishing and running university partnerships}: ``However, companies often struggle to establish and run university partnerships effectively.''
\end{itemize}

\paragraph{Personas.} The persona generator built three personas, each grounded in two observations:
\begin{itemize}
    \item \textbf{Linh}, 21, third-year marketing student who commutes and works part-time, grounded in the observations on matching by field and free time and on quick booking. Her goals are to find an advisor in marketing or branding whose hours fit around her classes and shifts, and to get a confirmed session in a few clicks.
    \item \textbf{Minh}, 29, senior software engineer and alumnus who mentors in his spare time, grounded in the observations on filtering requests and on a conflict-free schedule.
    \item \textbf{Huong}, 38, university relations manager at a technology company, grounded in the two observations on university partnerships.
\end{itemize}
The three personas carry different, partly conflicting priorities: the student wants fast booking, while the mentor wants to filter every request before accepting it.

\paragraph{Interview.} The interviews produced 22 messages. The first question to Linh asked what she did after two internship rejections, whom she turned to, and how long she waited for a useful answer. In her answer, the persona described turning first to her class group chat, where replies came quickly but were ``mostly sympathy'', and later explained that a mentor profile is only credible to her when it shows a specific title and company: ``If it just says `Marketing professional, 8 years experience' with no company --- skip.'' The plan had a topic on profile content and credential verification; the persona turned it into a concrete criterion of credibility, which the uploaded notes did not contain.

\paragraph{Checks of the evidence chain.} On the data of this run, every observation quote was found word for word in its source document (8 of 8), every observation named an existing source (8 of 8), every web source had a real URL (3 of 3), every evidence quote of a persona matched its observation exactly (6 of 6), and every persona used one of the stakeholder types of the plan (3 of 3).

\paragraph{Overview of the output.} Table~\ref{tab:impl-demo-output} summarizes the artifacts of the run. Overall, the prototype turned a one-sentence idea into a plan, a small grounded corpus, three distinct personas with traceable evidence and a first round of interviews, without manual intervention apart from the answers at the checkpoints. The output was coherent with the idea, the personas brought partly conflicting priorities into the interviews, and every link of the evidence chain could be checked. The run also showed the cost of the approach: it took about half an hour for the first half of the workflow, which confirms that cost and latency must be measured in the planned evaluation.

\begin{table}[H]
    \centering
    \caption{Overview of the output of the demonstration run.}
    \label{tab:impl-demo-output}
    \begin{tabularx}{\linewidth}{|l|X|}
        \hline
        \textbf{Artifact} & \textbf{Output} \\
        \hline
        Elicitation plan & Six stakeholder types and the topics of the interviews \\
        Research corpus & Three documents: the uploaded notes and two articles found by search \\
        Stakeholder observations & Eight observations, each with a verbatim quote from its document \\
        Personas & Three personas (a student, an alumnus mentor and an industry partner), each grounded in two observations \\
        Interview transcripts & Twenty-two messages \\
        Evidence checks & All checks passed: quotes (8 of 8), sources (8 of 8), URLs (3 of 3), persona evidence (6 of 6), stakeholder types (3 of 3) \\
        \hline
    \end{tabularx}
\end{table}
\section{Summary}
\label{sec:impl-summary}
This chapter described the implementation of \sysname{}: the technology stack and its rationale, the agent workflow on LangGraph with its state, versioned prompts and model layer, the backend with its background runs, event log, assistant and artifact services, and the web app generated from the workflow definition. More than seven hundred automated backend tests and the tests of the web app run in continuous integration on every change. A demonstration on a real idea showed that the prototype runs from the idea to persona interviews with real models and documents, and that its evidence chain can be checked. Chapter~\ref{chap:next-phase} plans how the framework will be evaluated and developed in the next phase.

